\documentclass[10pt,a4paper]{amsart}
\usepackage{amsmath,amssymb,amsthm,graphicx}
% Retained minimal preamble for the Stacks Project chapter more-algebra.tex (More on Algebra).
% The chapter file pinned-original-source.tex carries only \input{preamble} on its first line, so no
% preamble text lives inside the chapter itself.  The Stacks Project uses one single shared preamble for
% every chapter; the snapshot stages it identically as sources/stacks-chow/preamble.tex and
% sources/stacks-etale-cohomology/preamble.tex.  The block below is copied mechanically, line for line,
% from sources/stacks-chow/preamble.tex lines 14-18 (the Xy-pic setup) and 154-217 (the native theorem
% declarations together with the mathematical macros).  Nothing was added, removed or redefined; the
% wrapper class and the amsmath/amssymb/amsthm/graphicx base packages are supplied by the assembler.
% Chapter: more-algebra.tex -- More on Algebra.
% For commutative diagrams we use Xy-pic
\usepackage[all]{xy}

% We use 2cell for 2-commutative diagrams.
\xyoption{2cell}

% Theorem environments.
%
\theoremstyle{plain}
\newtheorem{theorem}[subsection]{Theorem}
\newtheorem{proposition}[subsection]{Proposition}
\newtheorem{lemma}[subsection]{Lemma}

\theoremstyle{definition}
\newtheorem{definition}[subsection]{Definition}
\newtheorem{example}[subsection]{Example}
\newtheorem{exercise}[subsection]{Exercise}
\newtheorem{situation}[subsection]{Situation}

\theoremstyle{remark}
\newtheorem{remark}[subsection]{Remark}
\newtheorem{remarks}[subsection]{Remarks}

\numberwithin{equation}{subsection}

% Macros
%
\def\lim{\mathop{\mathrm{lim}}\nolimits}
\def\colim{\mathop{\mathrm{colim}}\nolimits}
\def\Spec{\mathop{\mathrm{Spec}}}
\def\Hom{\mathop{\mathrm{Hom}}\nolimits}
\def\Ext{\mathop{\mathrm{Ext}}\nolimits}
\def\SheafHom{\mathop{\mathcal{H}\!\mathit{om}}\nolimits}
\def\SheafExt{\mathop{\mathcal{E}\!\mathit{xt}}\nolimits}
\def\Sch{\mathit{Sch}}
\def\Mor{\mathop{\mathrm{Mor}}\nolimits}
\def\Ob{\mathop{\mathrm{Ob}}\nolimits}
\def\Sh{\mathop{\mathit{Sh}}\nolimits}
\def\NL{\mathop{N\!L}\nolimits}
\def\CH{\mathop{\mathrm{CH}}\nolimits}
\def\proetale{{pro\text{-}\acute{e}tale}}
\def\etale{{\acute{e}tale}}
\def\QCoh{\mathit{QCoh}}
\def\Ker{\mathop{\mathrm{Ker}}}
\def\Im{\mathop{\mathrm{Im}}}
\def\Coker{\mathop{\mathrm{Coker}}}
\def\Coim{\mathop{\mathrm{Coim}}}

% Boxtimes
%
\DeclareMathSymbol{\boxtimes}{\mathbin}{AMSa}{"02}

%
% Macros for moduli stacks/spaces
%
\def\QCohstack{\mathcal{QC}\!\mathit{oh}}
\def\Cohstack{\mathcal{C}\!\mathit{oh}}
\def\Spacesstack{\mathcal{S}\!\mathit{paces}}
\def\Quotfunctor{\mathrm{Quot}}
\def\Hilbfunctor{\mathrm{Hilb}}
\def\Curvesstack{\mathcal{C}\!\mathit{urves}}
\def\Polarizedstack{\mathcal{P}\!\mathit{olarized}}
\def\Complexesstack{\mathcal{C}\!\mathit{omplexes}}
% \Pic is the operator that assigns to X its picard group, usage \Pic(X)
% \Picardstack_{X/B} denotes the Picard stack of X over B
% \Picardfunctor_{X/B} denotes the Picard functor of X over B
\def\Pic{\mathop{\mathrm{Pic}}\nolimits}
\def\Picardstack{\mathcal{P}\!\mathit{ic}}
\def\Picardfunctor{\mathrm{Pic}}
\def\Deformationcategory{\mathcal{D}\!\mathit{ef}}

% --- appended 2026-10-12 (Chief): the project own definitions of the reference/slogan/history
% display environments, copied mechanically from the shared Stacks preamble lines 8-12
% (\usepackage{verbatim} plus the three \newenvironment lines, which the chapter preamble keeps
% verbatim from the project preamble). Nothing added, removed or redefined: these environments
% carry no mathematical content and the project itself typesets them as comments. Without these
% three lines a counted statement whose own text uses them fails with Environment slogan undefined.
\usepackage{verbatim}
\newenvironment{reference}{\comment}{\endcomment}
\newenvironment{slogan}{\comment}{\endcomment}
\newenvironment{history}{\comment}{\endcomment}
\usepackage{cleveref}
\crefname{originallemma}{original lemma}{original lemmas}
\Crefname{originallemma}{Original lemma}{Original lemmas}
\makeatletter
\newlabel{algebra-lemma-completion-Noetherian-Noetherian}{{0316}{source}{}{sourcepointer-1}{}}
\newlabel{algebra-lemma-completion-Noetherian-Noetherian@cref}{{[originallemma][0][]0316}{[1][0][]source}{}{}{}}
\makeatother
\crefname{originallemma}{original lemma}{original lemmas}
\Crefname{originallemma}{Original lemma}{Original lemmas}
\makeatletter
\newlabel{algebra-lemma-completion-flat}{{00MB}{source}{}{sourcepointer-2}{}}
\newlabel{algebra-lemma-completion-flat@cref}{{[originallemma][0][]00MB}{[1][0][]source}{}{}{}}
\makeatother
\crefname{originallemma}{original lemma}{original lemmas}
\Crefname{originallemma}{Original lemma}{Original lemmas}
\makeatletter
\newlabel{algebra-lemma-ideal-is-squared-union-connected}{{00EH}{source}{}{sourcepointer-3}{}}
\newlabel{algebra-lemma-ideal-is-squared-union-connected@cref}{{[originallemma][0][]00EH}{[1][0][]source}{}{}{}}
\makeatother
\crefname{originallemma}{original lemma}{original lemmas}
\Crefname{originallemma}{Original lemma}{Original lemmas}
\makeatletter
\newlabel{lemma-characterize-henselian-pair}{{2553}{source}{}{sourcepointer-4}{}}
\newlabel{lemma-characterize-henselian-pair@cref}{{[originallemma][0][]2553}{[1][0][]source}{}{}{}}
\makeatother
\crefname{originallemma}{original lemma}{original lemmas}
\Crefname{originallemma}{Original lemma}{Original lemmas}
\makeatletter
\newlabel{lemma-complete-henselian}{{2482}{source}{}{sourcepointer-5}{}}
\newlabel{lemma-complete-henselian@cref}{{[originallemma][0][]2482}{[1][0][]source}{}{}{}}
\makeatother
\crefname{originallemma}{original lemma}{original lemmas}
\Crefname{originallemma}{Original lemma}{Original lemmas}
\makeatletter
\newlabel{lemma-glueing-sum-components-open}{{31984}{source}{}{sourcepointer-6}{}}
\newlabel{lemma-glueing-sum-components-open@cref}{{[originallemma][0][]31984}{[1][0][]source}{}{}{}}
\makeatother
\crefname{originallemma}{original lemma}{original lemmas}
\Crefname{originallemma}{Original lemma}{Original lemmas}
\makeatletter
\newlabel{lemma-lift-idempotent-upstairs}{{1978}{source}{}{sourcepointer-7}{}}
\newlabel{lemma-lift-idempotent-upstairs@cref}{{[originallemma][0][]1978}{[1][0][]source}{}{}{}}
\makeatother
\crefname{originalproposition}{original proposition}{original propositions}
\Crefname{originalproposition}{Original proposition}{Original propositions}
\makeatletter
\newlabel{proposition-equivalence}{{24813}{source}{}{sourcepointer-8}{}}
\newlabel{proposition-equivalence@cref}{{[originalproposition][0][]24813}{[1][0][]source}{}{}{}}
\makeatother
\crefname{originalremark}{original remark}{original remarks}
\Crefname{originalremark}{Original remark}{Original remarks}
\makeatletter
\newlabel{remark-formal-glueing-algebras}{{25036}{source}{}{sourcepointer-9}{}}
\newlabel{remark-formal-glueing-algebras@cref}{{[originalremark][0][]25036}{[1][0][]source}{}{}{}}
\makeatother
\crefname{originalremark}{original remark}{original remarks}
\Crefname{originalremark}{Original remark}{Original remarks}
\makeatletter
\newlabel{remark-glueing-data}{{24641}{source}{}{sourcepointer-10}{}}
\newlabel{remark-glueing-data@cref}{{[originalremark][0][]24641}{[1][0][]source}{}{}{}}
\makeatother
\title{A quotient of an adic completion killing a square-zero ideal of the kernel is again the completion of a finite type algebra}
\author{The Stacks Project Authors}
\date{stacks-project:more-algebra.tex:local-snapshot-20261012; GNU FDL 1.2 or later}
\begin{document}
\maketitle
\setcounter{subsection}{1}
\noindent\textit{Editorial note.} All delivered mathematical text is the author's own, copied from the pinned Stacks Project chapter file more-algebra.tex; no statement or proof step was rewritten, reordered or removed.\par\medskip
\noindent\textit{Editorial note.} The chapter file pinned-original-source.tex carries only \textbackslash{}input\{preamble\} on its first line, so no preamble text lives inside the chapter. The retained file sources/stacks-more-algebra/retained-preamble.tex is a mechanical, line-for-line copy of the project's single shared preamble as staged in the snapshot; nothing was added, removed or redefined and no package was added.\par\medskip
\noindent\textit{Editorial note.} The source URL is the source-level Stacks Project URL, as required when more than one card is produced from one source; the per-claim tag is carried by claim\_key.\par\medskip
\noindent\textit{Editorial note.} The author writes 'We omit the verification of conditions (1)(a) and (1)(b).' inside the counted proof span; that sentence is kept verbatim and marks the step left to the reader. The author's typo 'completetions' in the last paragraph is reproduced unchanged. The ten references inside the counted proof span are recorded as non-load-bearing navigation pointers; they mix cross-chapter tags with same-chapter label lines, and the arguments that consume them stay inside the counted proof span.\par\medskip
\noindent\textit{Allowed standard prerequisites (editorial).}\begin{itemize}
\item The gluing category Glue(B -> B-hat, f\_1, ..., f\_r) and the equivalence of the glueing theorem. invoked at proof lines 32056, 32070-32079. Is the formalism that converts the local splitting data into the finite B-algebra D.
\item Etale base change trivial modulo I and lifting of idempotents. invoked at proof lines 32084-32087. Produces the ring B' and the idempotent e that splits D.
\item Complete pairs are henselian, and the characterisation of henselian pairs. used at proof lines 32088-32091. Allows the map to C x C' to be factored through the etale extension, which is what makes the replacement of B legitimate.
\item Flatness of the completion map and the isomorphism B/IB -> B-hat/IB-hat. used at proof lines 32075-32077. Is the hypothesis that lets the glueing equivalence be applied to the completed situation.
\end{itemize}\medskip
\ifdefined\hypertarget\hypertarget{sourcepointer-1}{}\fi
\noindent\textit{Source reference.} Original lemma 0316: focused-seeds/AGP/sources/stacks-snapshot-20261011/tags.txt, source lines 3973--3973. Cross-chapter navigation pointer into the Stacks Project; the standard fact it supplies is applied inside the counted proof span and the argument resting on it stays there.\par\medskip
\ifdefined\hypertarget\hypertarget{sourcepointer-2}{}\fi
\noindent\textit{Source reference.} Original lemma 00MB: focused-seeds/AGP/sources/stacks-snapshot-20261011/tags.txt, source lines 817--817. Cross-chapter navigation pointer into the Stacks Project; the standard fact it supplies is applied inside the counted proof span and the argument resting on it stays there.\par\medskip
\ifdefined\hypertarget\hypertarget{sourcepointer-3}{}\fi
\noindent\textit{Source reference.} Original lemma 00EH: focused-seeds/AGP/sources/stacks-snapshot-20261011/tags.txt, source lines 534--534. Cross-chapter navigation pointer into the Stacks Project; the standard fact it supplies is applied inside the counted proof span and the argument resting on it stays there.\par\medskip
\ifdefined\hypertarget\hypertarget{sourcepointer-4}{}\fi
\noindent\textit{Source reference.} Original lemma 2553: sources/stacks-more-algebra/pinned-original-source.tex, source lines 2553--2553. Same-chapter navigation pointer; the statement it names is standard for the section and the argument that consumes it stays inside the counted proof span.\par\medskip
\ifdefined\hypertarget\hypertarget{sourcepointer-5}{}\fi
\noindent\textit{Source reference.} Original lemma 2482: sources/stacks-more-algebra/pinned-original-source.tex, source lines 2482--2482. Same-chapter navigation pointer; the statement it names is standard for the section and the argument that consumes it stays inside the counted proof span.\par\medskip
\ifdefined\hypertarget\hypertarget{sourcepointer-6}{}\fi
\noindent\textit{Source reference.} Original lemma 31984: sources/stacks-more-algebra/pinned-original-source.tex, source lines 31984--31984. Same-chapter navigation pointer; the statement it names is standard for the section and the argument that consumes it stays inside the counted proof span.\par\medskip
\ifdefined\hypertarget\hypertarget{sourcepointer-7}{}\fi
\noindent\textit{Source reference.} Original lemma 1978: sources/stacks-more-algebra/pinned-original-source.tex, source lines 1978--1978. Same-chapter navigation pointer; the statement it names is standard for the section and the argument that consumes it stays inside the counted proof span.\par\medskip
\ifdefined\hypertarget\hypertarget{sourcepointer-8}{}\fi
\noindent\textit{Source reference.} Original proposition 24813: sources/stacks-more-algebra/pinned-original-source.tex, source lines 24813--24813. Same-chapter navigation pointer; the statement it names is standard for the section and the argument that consumes it stays inside the counted proof span.\par\medskip
\ifdefined\hypertarget\hypertarget{sourcepointer-9}{}\fi
\noindent\textit{Source reference.} Original remark 25036: sources/stacks-more-algebra/pinned-original-source.tex, source lines 25036--25036. Same-chapter navigation pointer; the statement it names is standard for the section and the argument that consumes it stays inside the counted proof span.\par\medskip
\ifdefined\hypertarget\hypertarget{sourcepointer-10}{}\fi
\noindent\textit{Source reference.} Original remark 24641: sources/stacks-more-algebra/pinned-original-source.tex, source lines 24641--24641. Same-chapter navigation pointer; the statement it names is standard for the section and the argument that consumes it stays inside the counted proof span.\par\medskip
\begin{lemma}
\label{lemma-quotient-by-idempotent}
Let $A$ be a Noetherian ring and $I$ an ideal. Let $B$
be a finite type $A$-algebra. Let $B^\wedge \to C$ be a surjective
ring map with kernel $J$ where $B^\wedge$ is the $I$-adic completion.
If $J/J^2$ is annihilated by $I^c$ for some $c \geq 0$, then $C$ is
isomorphic to the completion of a finite type $A$-algebra.
\end{lemma}

\begin{proof}
Let $f \in I$. Since $B^\wedge$ is Noetherian (Algebra, Lemma
\ref{algebra-lemma-completion-Noetherian-Noetherian}),
we see that $J$ is a finitely generated ideal.
Hence we conclude from
Algebra, Lemma \ref{algebra-lemma-ideal-is-squared-union-connected}
that
$$
C_f = ((B^\wedge)_f)_e
$$
for some idempotent $e \in (B^\wedge)_f$. By
Lemma \ref{lemma-glueing-sum-components-open}
we can find a surjection $B^\wedge \to C'$
such that $B^\wedge \to C \times C'$ becomes an
isomorphism after inverting any $f \in I$.
Observe that $C \times C'$ is a finite $B^\wedge$-algebra.

\medskip\noindent
Choose generators $f_1, \ldots, f_r \in I$. Denote
$\alpha_i : (C \times C')_{f_i} \to B_{f_i} \otimes_B B^\wedge$
the inverse of the isomorphism of $(B^\wedge)_{f_i}$-algebras
we obtained above. Denote $\alpha_{ij} : (B_{f_i})_{f_j} \to (B_{f_j})_{f_i}$
the obvious $B$-algebra isomorphism.
Consider the object
$$
(C \times C', B_{f_i}, \alpha_i, \alpha_{ij})
$$
of the category $\text{Glue}(B \to B^\wedge, f_1, \ldots, f_r)$ introduced
in Remark \ref{remark-glueing-data}. We omit the verification of
conditions (1)(a) and (1)(b).
Since $B \to B^\wedge$ is a flat map
(Algebra, Lemma \ref{algebra-lemma-completion-flat}) inducing an isomorphism
$B/IB \to B^\wedge/IB^\wedge$ we may apply
Proposition \ref{proposition-equivalence}
and Remark \ref{remark-formal-glueing-algebras}.
We conclude that $C \times C'$ is isomorphic to $D \otimes_B B^\wedge$
for some finite $B$-algebra $D$.
Then $D/ID \cong C/IC \times C'/IC'$. Let $\overline{e} \in D/ID$
be the idempotent corresponding to the factor $C/IC$.
By Lemma \ref{lemma-lift-idempotent-upstairs} there exists an
\'etale ring map $B \to B'$ which induces an isomorphism
$B/IB \to B'/IB'$ such that $D' = D \otimes_B B'$ contains an
idempotent $e$ lifting $\overline{e}$. Since $C \times C'$
is $I$-adically complete the pair $(C \times C', IC \times IC')$
is henselian (Lemma \ref{lemma-complete-henselian}).
Thus we can factor the map $B \to C \times C'$ through $B'$ (see
Lemma \ref{lemma-characterize-henselian-pair}).
Doing so we may replace $B$ by $B'$ and $D$ by $D'$ (because the
$I$-adic completetions are unaltered by this). Then
we find that $D = D_e \times D_{1 - e} = D/(1 - e) \times D/(e)$
is a product of finite type $A$-algebras and the completion of the
first part is $C$ and the completion of the second part is $C'$.
\end{proof}

\end{document}
